\section*{Post-processing}
\begin{flushleft}

Post-processing of the Shor's Circuit results Fig. \ref{fig:qiskit_dist}, for a=4 and N=15, is the final step of Shor's Algorithm. From the ideal simulator, the measurement results were the states $|0000\rangle$ and $|1000\rangle$, with almost equal probability.

From equation (2), we know that the final state of the circuit will be:
\[
|2^m \phi\rangle
\]

Since we have 2 peaks, we can denote the final measurement as $l$, with:

\[
l_0 = [0000]_2 = [0]_{10}
\]

\[
l_1 = [1000]_2 = [1*2^4 + 0*2^3 + 0*2^1 + 0*2^0] = [8]_{10}
\]

and:
\[
l_i=2^m\phi_i, i \in \{0,1\}
\]
Therefore:

\[
\phi_0=0
\]
\[
\phi_1=8/2^4=1/2
\]

The period is encoded in the phase $\phi$, as:

\[
\phi = s/r
\]
\vspace{5cm}

Using continuous fractions to detect $r$

\[
\phi_0=\frac{s_0}{r_0} \xrightarrow{} \frac{0}{0}=\frac{s_0}{r_0}
\]

\[
\phi_1=\frac{s_1}{r_1} \xrightarrow{} \frac{1}{2}=\frac{s_1}{r_1}
\]

Therefore, we have two results:
\[
r_0=0, r_1=2
\]

Checking equation (1), to find the period:

The first solution of $r_0$ is not acceptable; thus we'll evaluate only $r_1$
\[
f(r_1) = a^{r_1} mod N = 4^2 mod 15 = 1
\]

Since f(2) = 1, we know that $r=2$.

Condition checking:
\[
4^{2/2} mod 15 = 4 \neq  \pm 1
\]

Factoring N, into two primes:
\[
p=gcd(r-1,N) = 3
\]

\[
q=gcd(r+1,N) = 5
\]
\end{flushleft}
